\scanpage{197}
penser que pour des pts différents à val. dans le \unsure{$=$|même} $k$, les deux foncteurs obtenus, soient isom (pas plus que dans le cas (i)). Il n’y a plus en général d’opérations de $\unsure{\Pi|\pi}_1(X)$ sur $\unsure{\xi|\zeta}_{x,\infty}$, mais encore comme dans (i) les autom de $k$ sur $k(x)$ opèrent dans $\xi_{x,\infty}$ $\tfrac12$ \uline{\unsure{linéaimt|linéairemt|linéairement}} (donnée de descente). Un autom$^{t}$ de $k$ qui n’est pas l’identité sur $k$ définit $\xi_{x,\infty}^{\sigma}$ qui à priori n’est plus isomorphe à $\xi_{x,\infty}$.

Un cas intéressant est celui où $k=\mathbb{Q}$, on trouve alors une \deleted{\unsure{trivialit.|trivial}} $\xi_{x\infty}\in T_{\unsure{\alpha|x}}(\mathbb{Q})$, i.e une trivialisation de $T_{\mathbb{Q}}$. Deux trivialisations à val dans $T_{\mathbb{Q}}$ ne sont en général pas isom à priori \unsure{les mêmes}. De plus, il n’y a pas d’isom \uline{canonique} prévisible entre

\[
\xi_{x,\infty}\ \text{et}\ \xi_{x\infty}\otimes_{\mathbb{Q}}\mathbb{Q}_{\ell}\ \text{et}\ \xi_{\bar{x},\ell}\otimes_{\mathbb{Z}_{\ell}}\mathbb{Q}_{\ell}\quad -
\]

\deleted{bien que pas} et comme le premier à priori ne provient pas d’un objet de $T_{\unsure{-|\alpha}}(\mathbb{Z}_\ell)$ mais seult de $T(\mathbb{Q}_\ell)$, il n’y a pas lieu de penser qu’il y ait $=$ un isom. Ainsi, tt pt rat / $\mathbb{Q}$ définit une classe de coh dans les $H^{2}(\mathbb{Q}_\ell,\unsure{\mathbb{G}|G}^{\circ(\alpha)}_{\ell})$. Dans le \deleted{\unsure{de}} cas général, il faut remplacer $\mathbb{Q}_\ell$

% slanted
\margin{en effet, ils sont déjà non isom pour les Motifs de poids et niveau $0$}

% archive sheet number: 196
